\subsection{The classical control computer and reversible computation}

\begin{definition}[Machines 1 and 2; uniform descriptions]\label{def:rtm}
Machine 1 is a deterministic multitape Turing machine with finite control
and alphabet and partial transition function, with designated stopping
states. It generates descriptions and handles genuinely observed classical
data. A family is uniform when one such machine on $1^n$ and classical
parameters outputs the complete finite instance and all its bounds.
For polynomial uniformity its output time and the numerical bounds are
polynomial in $n$ and the lengths of classical inputs. Resource parameters
used in compilation estimates are encoded in unary; wire addresses and
Boolean circuit descriptions are binary. A time bound supplied in binary (a horizon, in the sense of
Definitions~\ref{def:read} and \ref{def:adaptive}) still has the same
numerical execution cost, which may be exponential in its encoding length. Approximation requests are positive rational numbers.

Machine 2 is a deterministic reversible multitape Turing machine:
commands on each tape either read and replace a symbol or move the head
by $-1,0,1$ without reading or writing; distinct commands have disjoint
domains and disjoint ranges. Its global transition $\delta$ is partial
and injective. This is Bennett's quadruple convention
\cite[pp.~526--527, definition preceding the main theorem]{Bennett1973}.
For a finite encoding choose, on each bilateral tape, the window
$\{0,\ldots,S_0-1\}$, with $S_0\geq1$, and heads in that window.
For $k$ tapes, alphabet $\Gamma$, and control set $F$, use
\[
 m_{\rm RTM}=kS_0\lceil\log_2|\Gamma|\rceil+
 k\lceil\log_2 S_0\rceil+\lceil\log_2|F|\rceil
\]
bits, with padding if needed. Invalid codes and transitions leaving the
window are excluded from the domain $D$ of the restricted $\delta$.
On $\{0,1\}^{m_{\rm RTM}}\setminus D$ fix the lexicographic bijection
to the complement of $\delta(D)$; together with $\delta$ this defines
$\widehat\delta$. Its permutation matrix is the specified completion.
On excluded codes this is only a completion, not a fault-reporting TM
simulation. Any assertion of agreement with the original TM requires
every applied forward transition to remain in $D$. No such promise is
needed for unitarity or for applying the completed inverse.
The simulation in Definition~\ref{def:device} below uses explicit circuits
instead of this potentially expensive completion. There is no autonomous
cyclic scheduler in the present definition.
\end{definition}

\begin{lemma}[Finite completion]\label{lem:completion}
The completion $\widehat\delta$ is a permutation and hence defines a
unitary on its entire finite register.
\end{lemma}
\begin{proof}
 \step{r2comp-size}{$|D|=|\delta(D)|$.}
 \begin{proof}\pf\ The restriction is injective with image $\delta(D)$.
 \end{proof}
 \step{r2comp-complement}{The ordered complements have equal size and
 their prescribed correspondence is bijective.}
 \begin{proof}\pf\ Subtract the equality in \stepref{r2comp-size} from
 the size of the common ambient finite set.\end{proof}
 \qedstep
 \begin{proof}\pf\ The domains partition the set, and their images
 partition it into two disjoint sets, each mapped bijectively. A
 permutation takes an orthonormal basis to an orthonormal basis.\end{proof}
\end{proof}

\begin{definition}[Clean implementation workspace]\label{def:scratch}
An implementation of a logical unitary $A$ on $L$ is a total circuit
unitary $\widetilde A$ on $LZ$ satisfying
\[
 \widetilde A(\ket\psi\ket{0^{|Z|}})=
 A\ket\psi\ket{0^{|Z|}}\quad\hbox{for every }\ket\psi\in\cH_L.
\]
$Z$ is called clean workspace. This is a zero-input, zero-output
condition, not the assertion $\widetilde A=A\otimes\id_Z$ on other
workspace inputs. Logical recording cells are part of $L$, not $Z$.
For disjoint additional wires $Z'$, the padding of an operator $A$ is
$A\otimes\id_{Z'}$, with factors placed by their names. Padding a
clean implementation in this way preserves its defining equality.
\end{definition}

\begin{lemma}[Workspace and inverses]\label{lem:scratch}
Clean implementations compose, admit clean implementations of their
adjoints by circuit reversal, and may reuse the same zero workspace after
padding its width to the maximum required width.
\end{lemma}
\begin{proof}
 \step{r2scratch-invariant}{The subspace $\cH_L\otimes\ket0_Z$ is mapped
 onto itself by $\widetilde A$.}
 \begin{proof}\pf\ Its restriction is the unitary $A$ under the
 isometric identification $\psi\mapsto\psi\otimes\ket0$.\end{proof}
 \step{r2scratch-inverse}{$\widetilde A^\dagger(A\psi\otimes\ket0)
 =\psi\otimes\ket0$.}
 \begin{proof}\pf\ Apply $\widetilde A^\dagger$ to the defining equality;
 surjectivity of $A$ covers every logical input.\end{proof}
 \qedstep
 \begin{proof}\pf\ After each forward or inverse implementation the
 workspace is zero by the definition and \stepref{r2scratch-inverse}.
 Pad each workspace by tensor identity on the unused part of a common
 maximum-width pool. Its unused wires stay zero. Induction over the
 finite list proves composition and reuse.\end{proof}
\end{proof}

\begin{lemma}[Exact finite control and Boolean evaluation]\label{lem:synthesis}
NOT, CNOT, Toffoli, controlled-$H$, controlled-$T^{\pm1}$, and
controlled-CNOT have constant-size exact-phase implementations over
$\mathcal G^{\pm}$ with constant clean workspace. If a total Boolean
function $f:\{0,1\}^u\to\{0,1\}^v$ is supplied by an acyclic circuit
of $\beta$ AND, XOR, and NOT gates, then
$\ket{x,y}\mapsto\ket{x,y\oplus f(x)}$ has a clean implementation of
size and workspace $O(\beta+u+v)$. For a total Boolean predicate $p$ on
a register $X$ and a unitary $U$ on a disjoint target $R$, suppose
controlled-$U$ has a specified clean implementation. Computing $p(x)$
into a fresh flag, applying controlled-$U$, and undoing the predicate
implements $\sum_x\proj x_X\otimes U^{p(x)}$.
The controlled implementation acts on neither $X$ nor the predicate
evaluation workspace, apart from using the flag as an unchanged control.
\end{lemma}
\begin{proof}
 \step{r2synth-toffoli}{Toffoli has an exact implementation.}
 \begin{proof}
  \step{r2synth-phase}{For bits $a,b,c$,
  $4abc=a+b+c-(a\oplus b)-(a\oplus c)-(b\oplus c)
  +(a\oplus b\oplus c)$ as integers.}
  \begin{proof}\pf\ Substitute $a\oplus b=a+b-2ab$ and
  $a\oplus b\oplus c=a+b+c-2ab-2ac-2bc+4abc$.\end{proof}
  \step{r2synth-ccz}{There is a diagonal circuit with multiplier
  $(-1)^{abc}$ and zero-output workspace.}
  \begin{proof}\pf\ Compute the seven indicated parities into zero
  wires with CNOTs. Apply $T$ for positive coefficients and $T^\dagger$
  for negative coefficients, then erase the parities by reversing their
  CNOTs. The phase is $e^{i\pi(4abc)/4}$ by
  \stepref{r2synth-phase}.\end{proof}
  \qedstep
  \begin{proof}\pf\ Conjugate the third wire by $H$ before and after
  \stepref{r2synth-ccz}; $HZH=X$.\end{proof}
 \end{proof}
 \step{r2synth-controls}{The asserted controlled gates have exact phases.}
 \begin{proof}\pf\ Controlled-CNOT is Toffoli. For controlled-$T$,
 compute the AND of control and target into a zero wire, apply $T$ there,
 then uncompute; replace $T$ by $T^\dagger$ for its inverse. Put
 $V=SHTHS^\dagger$. Direct matrix multiplication gives $VZV^\dagger=H$.
 Thus $(\id\otimes V)\operatorname{CZ}(\id\otimes V^\dagger)$ is
 controlled-$H$, with $\operatorname{CZ}=(\id\otimes H)
 \operatorname{CNOT}(\id\otimes H)$. Scalar phases of $V$ cancel in
 this displayed expression. NOT is $HT^4H$.\end{proof}
 \step{r2synth-evaluate}{The Boolean evaluator has the asserted bounds.}
 \begin{proof}\pf\ Store each circuit gate's value in a fresh zero
 bit in topological order: AND by Toffoli, XOR by two CNOTs, NOT by
 CNOT followed by NOT on the new bit. Copy output bits into $y$, then
 undo the stored evaluations in reverse topological order. This uses
 $O(\beta+u+v)$ gates and bits
 and preserves $x$ on every logical input.\end{proof}
 \qedstep
 \begin{proof}\pf\ The construction for a predicate is the evaluator
 of \stepref{r2synth-evaluate}, the supplied controlled implementation,
 and the inverse evaluator. On $\ket x\ket\psi\ket0$ their successive
 logical values are $(x,\psi,p(x))$, $(x,U^{p(x)}\psi,p(x))$, and
 $(x,U^{p(x)}\psi,0)$. Linearity and Lemma~\ref{lem:scratch} give the
 asserted operator and zero workspace.\end{proof}
\end{proof}

\subsection{A total joint device tick}

A \emph{device} below is not a fourth machine. It is finite classical data,
supplied by Machine~1, that couples that control to Machine~3: at each
scheduled step (a \emph{tick}) the data select at most one quantum operation
and then update the classical configuration, with the previous configuration
retained so the step is a total unitary. Machine~2 enters only through its
completed permutation steps, admitted as atoms in Section~\ref{sec:instructions}.

\begin{definition}[Finite device data, not an implicit instruction table]
\label{def:device}
A register in the computational-basis state $\ket{0\cdots0}$ is called
\emph{blank}. A device instance supplies the following finite data. The quantum pool
$Q$ has $q$ named wires. Its classical configuration $C$ has $s$ bits,
including program position, word slots, memory, stack, quantum identifiers
and live flags, console tape and length, status, and outcome. Parameters
$w\geq8,B_{\rm mem},L_{\rm slot},h,o_{\max},p$ bound word width, memory bytes, word slots, call depth,
console bytes, and instruction positions respectively. Their field layout
and a validity predicate on all $s$-bit strings are part of the data.
A canonical valid fault configuration $c_{\rm fault}$ and an initial
configuration $c_{\rm in}$ are supplied.
The data also designate one active status code
$\sigma_{\mathrm{act}}$ and one timeout status code
$\sigma_{\mathrm{to}}$, each of the status-field width.

Crucially the data include total Boolean circuits for a decoder $d(c)$,
an update $f(c,z)$, and an output extraction $o(c)$. The word $z$ has a
fixed width $k\geq1$, the maximum read width. Write $[q]=\{0,\ldots,q-1\}$.
The decoder's semantic codomain is the disjoint union
\[
 \begin{split}
 \mathcal D={}&\{\mathsf{ordinary},\mathsf{idle},\mathsf{fault}\}
 \sqcup\{\mathsf{gate}(g,i,j):g\in\{H,T,\operatorname{CNOT}\}\}\\
 &\sqcup\{\mathsf{read}(r;i_0,\ldots,i_{r-1}):0\leq r\leq k\}
 \sqcup\{\mathsf{measure}(i),\mathsf{allocate}(i):i\in[q]\}.
 \end{split}
\]
In gate and read modes all used addresses belong to $[q]$; CNOT has
$i\ne j$; read lists may repeat addresses. A one-wire gate uses
only $i$. The fixed binary layout, in order, is a 3-bit mode tag (the
seven displayed modes in order, starting at zero), a 2-bit gate tag
($H,T,\operatorname{CNOT}$ numbered $0,1,2$), two address fields, a
read-length field, and $k$ read-address fields. Measurement and allocation
use the first address field; one-wire gates do likewise. Address width is
$a_q=\max(1,\lceil\log_2(q+1)\rceil)$; length width is
$a_k=\lceil\log_2(k+1)\rceil$. Unused fields are zero. Thus decoder
width is $5+(k+2)a_q+a_k$. No alternative redundant encoding is admitted.
The supplied Boolean decoder always outputs a canonical code in
$\mathcal D$; a malformed intended action must instead encode fault.
A read XORs wire $i_j$ into sample bit $j$, starting at bit zero;
measurement XORs wire $i$ into sample bit zero only. The sample's other
bits are untouched, hence zero on a blank sample. Sample destinations
in $C$ are determined by $f$ alone, not by additional decoder fields.

For every configuration code, a gate, read, or measurement may address
only wires whose live flags are one. Allocation selects a zero live
flag and sets just that flag to one; its quantum action is identity on
that preallocated factor. Except when sending to $c_{\rm fault}$, all
other updates preserve live flags for every $z$; allocation preserves
the other flags. There is no deallocation instruction. Invalid codes,
addresses, repeated gate operands, exhausted capacities, mode violations,
or illegal instructions have decoder fault, identity quantum action,
and $f(c,z)=c_{\rm fault}$ for every $z$. The fault configuration is valid
and terminal. These are requirements on the total supplied circuits,
not an implicit run-time repair. In particular no gate may precede
allocation of its operand in a valid forward execution.
Valid terminal configurations have idle decoder and $f(c,z)=c$.
For ordinary, gate, allocation and idle modes, $f$ is independent of $z$.
For read and measurement it writes the sampled value into its specified
destination and performs the specified program-position update.

All remaining behavior is supplied by the circuits, including which
operands are immediate or indirect, signed comparison, shift directions,
division, call-frame layout, console overflow, and every output bit on
invalid encodings. No theorem identifies unspecified choices.
For example modular addition means addition modulo $2^w$ only if that
equation is imposed on the supplied circuit. This parametrization fixes
one operator for every complete instance. Instruction names 1--6 mean
the corresponding classical functions when the instance implements them;
no efficiency of an arbitrary supplied function is presumed.
The circuit for $f$ includes instruction selection and has size $\beta$.
Let $\Delta$ bound the sum of sizes of decoder and output circuits,
including their input and output ports. The encoding has
$s\geq a_q$; in particular the live-flag field contains $q$ bits.
\end{definition}

For a first reading, keep one small schematic instance in mind. Take a
pool of two wires, a configuration $C$ holding a two-bit program position,
two live flags and a status bit, with the other required fields padded to
the device layout. Total circuits $d,f$ decode position $0$ as $H$ on
wire $0$ and advance to position $1$, decode position $1$ as halted/idle,
and send every other position code to the fault configuration, subject to
the validity and live-flag checks above. Nothing in what follows depends on
the size of the fields: a device is simply a finite classical control
state $C$ together with total Boolean circuits saying, for \emph{every}
value of $C$, which single quantum operation (if any) happens next and
what $C$ becomes. The binary layout and the every-code conditions exist
so that this statement has no undefined cases when $C$ itself is later
held in qubits.

One classical update $c\mapsto f(c,z)$ is in general many-to-one, so it
has no inverse and no unitary. The construction below makes it reversible
by retaining the old configuration in a fresh register $J_t$ before
overwriting $C$. The map
$(c,j,z)\mapsto(j\oplus f(c,z),c,z)$ is then a permutation of \emph{all}
basis states, blank or not, with the displayed inverse; that totality is
what later lets a rewind act after a keep has moved the state outside the
image of the forward steps on blank records. The price is the distinction
used throughout: the operator is a unitary unconditionally, but it means
``perform the classical update'' only on blank $J_t,B_t$, where the full
tick acts as $(c,0,0)\mapsto(f(c,z),c,z)$ on the machine factors.
The order in $S_t=F_tD_t$ is essential: the quantum operation $D_t$ is
selected by the configuration \emph{before} the tick, and only afterwards
does $F_t$ change that configuration.

\begin{definition}[Snapshot, sampling, and ordering]\label{def:joint-tick}
For each scheduled tick $t$ reserve distinct registers $J_t$ of $s$ bits
and $B_t$ of $k$ bits, initially zero. They are retained. Define
\[
 F_t:\ket{c,j,z}\longmapsto
       \ket{j\oplus f(c,z),c,z}_{C J_t B_t},\qquad
 F_t^{-1}:\ket{c',j',z}\longmapsto
       \ket{j',c'\oplus f(j',z),z}.
\]
Define $G_c$ on $QB_t$ using the \emph{pre-tick} decoder: the selected
gate in gate mode; CNOTs from selected pool wires to the designated
positions of $B_t$ in read or measurement mode; identity in all other
modes. In particular coherent read and measurement have the same retained
copying operation, but only the latter has an observational interpretation.
The joint tick on the full retained list is
\[
 D_t=\sum_{c\in\{0,1\}^s}\proj c_C\otimes G_c,
 \qquad S_t=F_tD_t.                         \tag{1}\label{eq:tick}
\]
No decoder field changes during $D_t$. The update follows the quantum
action, so an address changed by $f$ cannot change that action's operand.
The logical machine register $M$ includes $C$, every $J_t,B_t$, and any
other explicit retained records; implementation workspace is separate.
For an ordinary tick with blank records,
$(c,0,0)\mapsto(f(c,0),c,0)$. For a read of basis value $x$ it is
$(c,0,0)\mapsto(f(c,x),c,x)$, with the read pool wires unchanged.

All syntactic tick occurrences of a complete term receive globally
distinct $J_t,B_t$ names, including occurrences in sequential fragments
and different marked bodies. Expanding a term's adjoint uses the same
names as the term, in inverse order; it allocates no new record.
The ordinary interpretation of a forward tick requires its records blank
at that application. If a nonblank cell is encountered the completed
unitary still applies, but it is not asserted to implement the ordinary
update. Fresh cells may be reused only with a proved zero-restoration
condition. Every name possibly read by an enclosing inverse is retained.
\end{definition}

\begin{lemma}[Joint tick well-posedness]\label{lem:joint}
Equation~\eqref{eq:tick} is a total unitary with inverse
$D_t^\dagger F_t^{-1}$, independent of blank-record promises. It has a
clean implementation over $\mathcal G^{\pm}$. On blank records it has
the stated ordinary and coherent-read actions.
\end{lemma}
\begin{proof}
 \step{r2joint-permutation}{$F_t^{-1}F_t=F_tF_t^{-1}=\id$.}
 \begin{proof}\pf\ Substitute the displayed binary formulas and use
 $a\oplus a=0$, while $z$ is preserved.\end{proof}
 \step{r2joint-controlled}{$D_t$ is unitary.}
 \begin{proof}\pf\ Each $G_c$ is a product of unitary gates and the
 projectors $\proj c$ are orthogonal and sum to identity. Hence both
 $D_t^\dagger D_t$ and $D_tD_t^\dagger$ equal identity.\end{proof}
 \step{r2joint-implementation}{Both factors have clean implementations.}
 \begin{proof}\pf\ Evaluate $f(c,z)$ by Lemma~\ref{lem:synthesis}, XOR
 it into $J_t$, restore workspace, and swap $C,J_t$ using three CNOTs
 per pair. For $D_t$, compute the decoder into workspace, enumerate
 gate/measurement actions and read position/address pairs, compute
 their selection predicates, apply each controlled gate, and uncompute.
 For reads, the predicate also tests that the position is below the
 decoded length. The decoder depends only on the unchanged $C$;
 selected read operations are controlled-CNOTs. Uncompute the decoder.
 Lemmas~\ref{lem:synthesis} and \ref{lem:scratch} apply.\end{proof}
 \qedstep
 \begin{proof}\pf\ Multiply the inverses from
 \stepref{r2joint-permutation} and \stepref{r2joint-controlled} in reverse
 order. Substitution of $j=z=0$ gives the stated forward actions.
 \stepref{r2joint-implementation} supplies the realization.\end{proof}
\end{proof}

